\documentclass[12pt,a4paper]{article}
\usepackage[utf8]{inputenc}
\usepackage[italian]{babel}
\usepackage{amsmath, amssymb, bm}
\usepackage{graphicx}
\usepackage{pgfplots}
\pgfplotsset{compat=1.18}
\usepackage{geometry}
\geometry{margin=2.5cm}
\usepackage{caption}
\usepackage{hyperref}
\usepackage{tikz}
\usetikzlibrary{arrows.meta, positioning, shapes}

\title{\centering Modello integrato della Coscienza Umana \\ Teoria Quadristato}
\author{Autore: Giuseppe Carlino}
\date{\today}

\begin{document}
\maketitle

\section*{Schema concettuale integrato}

\begin{center}
\begin{tikzpicture}[node distance=2.5cm, every node/.style={align=center, font=\small}]
    % Nodi principali
    \node[draw, rectangle, fill=blue!20] (tempo) {Legge 1\\Tempo percepito $t_c$};
    \node[draw, rectangle, fill=green!20, right=of tempo] (conoscenza) {Legge 0\\Conoscenza $k$};
    \node[draw, rectangle, fill=orange!20, below=of conoscenza] (scelta) {Legge 2\\Scelta $s$};
    \node[draw, rectangle, fill=purple!20, left=of scelta] (volonta) {Legge 3\\Volontà $w$};

    % Frecce principali (interazioni)
    \draw[->, thick] (tempo) -- (conoscenza) node[midway, above] {influenza crescita};
    \draw[->, thick] (conoscenza) -- (scelta) node[midway, right] {guida decisioni};
    \draw[->, thick] (scelta) -- (volonta) node[midway, below] {trasforma in azione};
    \draw[->, thick] (volonta) -- (scelta) node[midway, left] {retroazione};

    % Retroazione tempo
    \draw[->, thick, dashed] (conoscenza.north) .. controls +(up:2cm) and +(up:2cm) .. (tempo.north) 
        node[midway, above] {influenza percezione};

    % Perturbazioni
    \node[draw, ellipse, fill=red!20, above=of tempo] (perturbazioni) {Perturbazioni $\eta$};
    \draw[->, thick, dotted] (perturbazioni) -- (tempo);
    \draw[->, thick, dotted] (perturbazioni) -- (conoscenza);
    \draw[->, thick, dotted] (perturbazioni) -- (scelta);
    \draw[->, thick, dotted] (perturbazioni) -- (volonta);

    % Legge generale (quadristato)
    \node[draw, rounded rectangle, thick, fill=yellow!20, below=of volonta, minimum width=10cm] (quadristato) {Legge 4\\Coscienza Quadristato\\$\mathbf{X}=(t_c, k, s, w)$};
    \draw[->, thick, dashed] (tempo.south) -- (quadristato.north);
    \draw[->, thick, dashed] (conoscenza.south) -- (quadristato.north);
    \draw[->, thick, dashed] (scelta.south) -- (quadristato.north);
    \draw[->, thick, dashed] (volonta.south) -- (quadristato.north);
\end{tikzpicture}
\end{center}

\section*{Legenda concettuale}
\begin{itemize}
    \item Rettangoli colorati: leggi individuali della coscienza.
    \item Frecce continue: flussi principali e retroazioni tra stati.
    \item Frecce tratteggiate: retroazione globale verso il modello quadridimensionale.
    \item Frecce puntinate rosse: perturbazioni interne ed esterne che influenzano ogni stato.
    \item Rettangolo giallo: rappresenta la Legge Generale (Quadristato), integrando tutti e quattro gli stati.
\end{itemize}

\end{document}